\section{Related Work}
\label{sec:related}

\textbf{Ticket triage stops at the label.} Montgomery et al.\
\cite{montgomery2018escalation} model escalation risk over 2.5M IBM support
tickets; DeepTriage \cite{mani2019deeptriage} assigns bug reports to developers;
Lee et al.\ \cite{lee2017industrialtriage} report an industrial incident
deployment. All three predict a per-ticket label and report classification
metrics, and none asks what the prediction is worth once tickets compete for a
finite number of agents. That question is this paper's subject.

\textbf{Scheduling answers the ordering question.} The $c\mu$ rule
\cite{cox1961queues} is the classical result for linear delay cost, with
Cobham's priority analysis \cite{cobham1954priority} and Kleinrock's
delay-dependent priorities \cite{kleinrock1964delay} preceding and extending it,
accumulating-priority queues adding aging \cite{stanford2014accumulating}, and
the generalised $c\mu$ rule \cite{vanmieghem1995generalized} covering convex
cost. Our score is best understood as a $c\mu$ rule whose cost term is
\emph{learned} rather than assigned: the classical rule assumes a known per-class
cost, and ours estimates it per ticket from a posterior. Argon and Ziya
\cite{argon2009priority} give the result that licenses this directly, namely that
ordering by a posterior index rather than by its argmax lowers long-run cost,
with the gain growing in the variability of the index. We use
\cite{pinedo2022scheduling} for objectives and nDCG \cite{jarvelin2002ndcg} for
queue-head quality.

\textbf{Fairness here is a scheduling property.} Our disparity result concerns
groups of \emph{customers} and is denominated in waiting time, which
distinguishes it from fairness in ranking \cite{singh2018fairexposure}, where
the groups are ranked items and the currency is exposure.

\textbf{Multilingual and romanized input.} The classifiers feeding our queue are
multilingual encoders \cite{conneau2020xlmr,feng2022labse,marone2025mmbert} and
Indic-targeted models \cite{khanuja2021muril}, evaluated on Sinhala and Tamil
written in both native and Latin script. Human-typed romanized Tamil
\cite{chakravarthi2020tamilmix} is the reference point that makes our
synthetically romanized tracks an optimistic bound; Sinhala resources remain
thin \cite{desilva2019sinhalasurvey,sumanathilaka2026swabhasha}. Intent labels
come from BANKING77 \cite{casanueva2020banking77}. Because two of our three label
columns are model-generated, we report measured ceilings throughout, following
the treatment of benchmark label error in \cite{northcutt2021labelerrors}.
